\appendix
\chapter{Source Code Listings}
\label{app:source_code}

This appendix contains the complete source code for the sleep stage classification
system. Python listings cover the data processing, feature engineering, model
training, and evaluation pipeline. Arduino/C listings cover the nRF52840 firmware.
All Python code was developed under Python~3.12 and tested on macOS Sequoia.
Arduino firmware was compiled and flashed using Arduino IDE~2.x with the Seeed
nRF52 board package.

\section{Environment Setup}
\label{app:env_setup}

\begin{bashcode}
# Create and activate virtual environment
python3 -m venv venv
source venv/bin/activate        # macOS / Linux

# Install all dependencies
pip install -r requirements.txt
\end{bashcode}

\newpage
\section{Repository Structure and Reproducibility Seed}
\label{app:repo_structure}

\begin{verbatim}

Sleep_Stage_Classifier/
├── data/                     # Raw data (excluded from Git)
├── src/
│   ├── preprocessing/        # Data synchronisation
│   ├── feature_engineering/  # Feature extraction
│   ├── modeling/             # Model training
│   └── evaluation/           # Performance analysis
├── results/                  # Figures and CSV results
├── hardware/
│   ├── model_export/         # micromlgen export scripts
│   └── firmware/
│       └── SleepClassifier_Audio_XGBoost/   # nRF52840 Arduino firmware
├── requirements.txt
├── .gitignore
└── README.md

\end{verbatim}

\begin{pycode}
import numpy as np
import random

SEED = 42
np.random.seed(SEED)
random.seed(SEED)
\end{pycode}

\section{Data Loading}
\label{app:data_loading}

\begin{pycode}
import pandas as pd
import pyedflib

def load_psg_data(edf_path, annotation_path):
    """
    Load PSG sleep stage annotations from CSV/XML.
    Returns DataFrame with 3-class consolidated labels.
    """
    annotations = pd.read_csv(annotation_path)

    stage_mapping = {
        'Wake': 0, 'W': 0,
        'N1': 1, 'N2': 1, 'N3': 1, 'N4': 1,
        'REM': 2, 'R': 2
    }
    annotations['sleep_stage_3class'] = (
        annotations['sleep_stage'].map(stage_mapping)
    )
    return annotations
\end{pycode}

\begin{pycode}
def load_imu_data(csv_path):
    """
    Load IMU sensor data from PillowClip device CSV.
    Columns: ts_ms, ax, ay, az, gx, gy, gz, tempC
    """
    imu_data = pd.read_csv(csv_path)
    imu_data['timestamp'] = imu_data['ts_ms'] / 1000.0
    imu_data = resample_to_1hz(imu_data)
    return imu_data
\end{pycode}

\section{Data Synchronisation}
\label{app:synchronisation}

\begin{pycode}
def synchronize_timestamp_based(psg_annotations, imu_data):
    """
    Align PSG and IMU using absolute timestamps.
    Requires both systems to have matching absolute time references.
    """
    synchronized_epochs = []

    for idx, row in psg_annotations.iterrows():
        epoch_start = row['timestamp']
        epoch_end   = epoch_start + 30

        mask = ((imu_data['timestamp'] >= epoch_start) &
                (imu_data['timestamp'] <  epoch_end))
        imu_epoch = imu_data[mask]

        if len(imu_epoch) >= 25:   # 5-sample tolerance
            synchronized_epochs.append({
                'patient_id':  row['patient_id'],
                'epoch_id':    idx,
                'sleep_stage': row['sleep_stage_3class'],
                'imu_data':    imu_epoch
            })

    return synchronized_epochs
\end{pycode}

\begin{pycode}
def synchronize_sequential(psg_annotations, imu_data,
                            samples_per_epoch=30):
    """
    Align PSG and IMU via sequential epoch matching.
    Used when precise start-time metadata is unavailable.
    """
    synchronized_epochs = []
    n_epochs = len(psg_annotations)

    for i in range(n_epochs):
        start_idx = i * samples_per_epoch
        end_idx   = start_idx + samples_per_epoch

        if end_idx <= len(imu_data):
            imu_epoch = imu_data.iloc[start_idx:end_idx]
            synchronized_epochs.append({
                'patient_id':  psg_annotations.iloc[i]['patient_id'],
                'epoch_id':    i,
                'sleep_stage': psg_annotations.iloc[i]['sleep_stage_3class'],
                'imu_data':    imu_epoch
            })

    return synchronized_epochs
\end{pycode}

\begin{pycode}
def validate_synchronization(synchronized_epochs):
    """
    Validate synchronisation quality via wake/sleep variance ratio.
    Accepts sessions with ratio > 1.5.
    """
    wake_var  = []
    sleep_var = []

    for epoch in synchronized_epochs:
        imu      = epoch['imu_data']
        variance = np.var(imu['ax'])
        if epoch['sleep_stage'] == 0:
            wake_var.append(variance)
        else:
            sleep_var.append(variance)

    ratio = np.mean(wake_var) / (np.mean(sleep_var) + 1e-6)
    print(f"Wake/Sleep variance ratio: {ratio:.2f}  (threshold: >1.5)")
    return ratio > 1.5
\end{pycode}


\section{Feature Extraction}
\label{app:feature_extraction}

\begin{pycode}
def extract_all_features(epoch_data, patient_id, patient_stats):
    """
    Top-level dispatcher: extracts all 98 IMU features for one epoch.
    The final 118-feature model appends 20 audio features separately.

    Args:
        epoch_data    : DataFrame, 30 rows x 6 IMU axes
        patient_id    : str, participant identifier
        patient_stats : dict, pre-computed per-patient baselines

    Returns:
        features : dict, 98 key-value pairs
    """
    features = {}
    features.update(extract_baseline_features(epoch_data))        # 36
    features.update(extract_movement_features(epoch_data))        # 12
    features.update(extract_patient_features(                     # 14
        epoch_data, patient_id, patient_stats))
    features.update(extract_temporal_features(epoch_data))        # 15
    features.update(extract_orientation_features(epoch_data))     #  7
    features.update(extract_variability_features(epoch_data))     #  6
    features.update(extract_bout_features(epoch_data))            #  5
    features.update(extract_circadian_features(epoch_data))       #  3
    return features
\end{pycode}

\begin{pycode}
def extract_movement_features(epoch_data):
    """Orientation-invariant magnitude features (12 features)."""
    features = {}

    ax, ay, az = epoch_data['ax'], epoch_data['ay'], epoch_data['az']
    acc_mag    = np.sqrt(ax**2 + ay**2 + az**2)

    features['acc_mag']     = np.mean(acc_mag)
    features['acc_mag_std'] = np.std(acc_mag)
    features['acc_mag_max'] = np.max(acc_mag)
    features['acc_mag_min'] = np.min(acc_mag)

    gx, gy, gz = epoch_data['gx'], epoch_data['gy'], epoch_data['gz']
    gyro_mag   = np.sqrt(gx**2 + gy**2 + gz**2)

    features['gyro_mag']     = np.mean(gyro_mag)
    features['gyro_mag_std'] = np.std(gyro_mag)
    features['total_movement'] = (features['acc_mag'] +
                                  features['gyro_mag'])

    threshold = 0.5   # g
    features['movement_bout_count'] = count_threshold_crossings(
        acc_mag, threshold)

    stillness_threshold = 0.2   # g
    features['stillness_duration'] = count_consecutive_below_threshold(
        acc_mag, stillness_threshold)

    return features
\end{pycode}

\begin{pycode}
def extract_patient_features(epoch_data, patient_id, patient_stats):
    """Patient-specific normalisation features (14 features)."""
    features = {}
    stats = patient_stats[patient_id]

    acc_mag = np.mean(np.sqrt(
        epoch_data['ax']**2 +
        epoch_data['ay']**2 +
        epoch_data['az']**2
    ))

    features['patient_acc_mean']   = stats['acc_mean']
    features['patient_acc_median'] = stats['acc_median']
    features['patient_acc_std']    = stats['acc_std']
    features['patient_wake_ratio'] = stats['wake_ratio']

    features['acc_mag_normalized'] = (
        (acc_mag - stats['acc_mean']) / (stats['acc_std'] + 1e-6)
    )
    # relative_activity is an IMU patient-normalisation feature:
    # current epoch magnitude relative to patient's baseline median.
    features['relative_activity'] = (
        acc_mag / (stats['acc_median'] + 1e-6)
    )
    features['acc_percentile'] = compute_percentile_rank(
        acc_mag, stats['acc_distribution']
    )
    return features
\end{pycode}


\section{Dataset Serialisation}
\label{app:serialisation}

\begin{pycode}
import pickle
import numpy as np

def create_final_dataset(synchronized_epochs, patient_stats,
                          output_path='models/sleep_dataset_optimized.pkl'):
    """
    Runs feature extraction over all epochs and serialises the result.
    """
    X, y, patient_ids = [], [], []

    for epoch in synchronized_epochs:
        features = extract_all_features(
            epoch['imu_data'], epoch['patient_id'], patient_stats
        )
        X.append(list(features.values()))
        y.append(epoch['sleep_stage'])
        patient_ids.append(epoch['patient_id'])

    X           = np.array(X)
    y           = np.array(y)
    patient_ids = np.array(patient_ids)

    dataset = {
        'X_all_features': X,
        'y_3class':        y,
        'patient_ids':     patient_ids,
        'feature_names':   list(features.keys()),
        'n_features':      X.shape[1],
        'n_samples':       X.shape[0],
        'class_distribution': {
            'Wake': int(np.sum(y == 0)),
            'NREM': int(np.sum(y == 1)),
            'REM':  int(np.sum(y == 2)),
        }
    }

    with open(output_path, 'wb') as f:
        pickle.dump(dataset, f)

    print(f"Saved: {dataset['n_samples']} samples, "
          f"{dataset['n_features']} features -> {output_path}")
    return dataset
\end{pycode}


\section{Model Training}
\label{app:training}

\begin{pycode}
from sklearn.model_selection import GroupKFold
from sklearn.preprocessing import StandardScaler
from sklearn.utils.class_weight import compute_class_weight
from sklearn.metrics import cohen_kappa_score, accuracy_score, f1_score
from xgboost import XGBClassifier
import pickle, numpy as np

def train_xgboost_model(dataset_path, n_splits=5):
    """
    Trains XGBoost with 5-fold patient-level GroupKFold CV and
    class-balanced sample weights.  The cross-validation study uses
    n_estimators=100; the final embedded model uses n_estimators=20
    to satisfy the nRF52840 1 MB flash constraint.
    """
    with open(dataset_path, 'rb') as f:
        data = pickle.load(f)

    X, y, groups = (data['X_all_features'],
                    data['y_3class'],
                    data['patient_ids'])
    gkf = GroupKFold(n_splits=n_splits)

    kappa_scores, acc_scores, f1_scores = [], [], []

    for fold, (train_idx, test_idx) in enumerate(
            gkf.split(X, y, groups)):
        print(f"\n=== Fold {fold + 1}/{n_splits} ===")

        scaler  = StandardScaler()
        X_train = scaler.fit_transform(X[train_idx])
        X_test  = scaler.transform(X[test_idx])
        y_train, y_test = y[train_idx], y[test_idx]

        classes = np.unique(y_train)
        cw = compute_class_weight('balanced',
                                  classes=classes, y=y_train)
        sample_weights = np.array(
            [dict(zip(classes, cw))[yi] for yi in y_train]
        )

        model = XGBClassifier(
            n_estimators=100,
            max_depth=6,
            learning_rate=0.1,
            random_state=42,
            eval_metric='mlogloss',
            verbosity=0
        )
        model.fit(X_train, y_train, sample_weight=sample_weights)
        y_pred = model.predict(X_test)

        kappa = cohen_kappa_score(y_test, y_pred)
        acc   = accuracy_score(y_test, y_pred)
        f1    = f1_score(y_test, y_pred, average='macro')

        kappa_scores.append(kappa)
        acc_scores.append(acc)
        f1_scores.append(f1)

        print(f"Kappa: {kappa:.3f}  Accuracy: {acc:.3f}  "
              f"F1-Macro: {f1:.3f}")

    print(f"\n=== Mean Results ===")
    print(f"Kappa:    {np.mean(kappa_scores):.3f} "
          f"+- {np.std(kappa_scores):.3f}")
    print(f"Accuracy: {np.mean(acc_scores):.3f} "
          f"+- {np.std(acc_scores):.3f}")
    print(f"F1-Macro: {np.mean(f1_scores):.3f} "
          f"+- {np.std(f1_scores):.3f}")

    return model
\end{pycode}


\section{Class Weight Computation}
\label{app:class_weights}

\begin{pycode}
from sklearn.utils.class_weight import compute_class_weight
import numpy as np

def get_sample_weights(y_train):
    """
    Computes per-sample balanced class weights using scikit-learn.
    Equivalent to w_c = N / (k * N_c) for each class c.

    Returns a weight array aligned with y_train for use as
    sample_weight in XGBClassifier.fit().
    """
    classes = np.unique(y_train)
    cw = compute_class_weight('balanced',
                               classes=classes, y=y_train)
    weight_map = dict(zip(classes, cw))

    sample_weights = np.array([weight_map[yi] for yi in y_train])
    print(f"Class weights: { {int(k): round(v, 3)
                               for k, v in weight_map.items()} }")
    return sample_weights
\end{pycode}


\section{Feature Importance Analysis}
\label{app:feature_importance}

\begin{pycode}
import pandas as pd

def extract_feature_importance(model, feature_names,
                                output_csv='results/feature_importance.csv'):
    """
    Extracts XGBoost Gini importance, ranks features, and saves to CSV.
    """
    importance_df = pd.DataFrame({
        'feature':    feature_names,
        'importance': model.feature_importances_
    }).sort_values('importance', ascending=False)

    importance_df.to_csv(output_csv, index=False)

    print("\nTop 10 Features:")
    print(importance_df.head(10).to_string(index=False))

    return importance_df
\end{pycode}


\section{Model Export to C Header}
\label{app:model_export}

\begin{pycode}
from micromlgen import port
import pickle

def export_xgboost_to_c(
        model_pkl_path,
        output_path='hardware/firmware/'
                    'SleepClassifier_Audio_XGBoost/'
                    'sleep_model_audio/xgboost_audio.h'):
    """
    Exports a trained XGBoost model to a C header for nRF52840 deployment
    using the micromlgen library, which generates an
    Eloquent::ML::Port::XGBClassifier class directly callable from Arduino.

    Note: micromlgen requires the full sklearn XGBClassifier object
    (not the raw booster JSON).
    """
    with open(model_pkl_path, 'rb') as f:
        model = pickle.load(f)

    c_code = port(model)

    with open(output_path, 'w') as f:
        f.write(c_code)

    print(f"Model exported -> {output_path}")
\end{pycode}


\section{Arduino Firmware --- Main Loop}
\label{app:firmware}

\begin{ccode}
// SleepClassifier_Audio_XGBoost.ino
// nRF52840 main loop: 1 Hz IMU sampling, PDM microphone streaming,
// 30-second epoch classification, and RGB LED sleep stage feedback.

#include <Wire.h>
#include <LSM6DS3.h>
#include <PDM.h>
#include "sleep_model_audio/xgboost_audio.h"   // micromlgen-generated
#include "sleep_model_audio/scaler_audio.h"     // StandardScaler params
#include "sleep_model_audio/features_audio.h"   // FEAT_* index defines

LSM6DS3 imu(I2C_MODE, 0x6A);

float imu_buffer[30][6];   // 30 samples x [ax, ay, az, gx, gy, gz]
int   sample_count = 0;

void setup() {
    Serial.begin(115200);

    if (imu.begin() != 0) {
        Serial.println("IMU init failed!");
        while (1);
    }

    // PDM microphone (16 kHz, mono)
    PDM.onReceive(on_pdm_data);
    if (!PDM.begin(1, 16000)) {
        Serial.println("PDM init failed!");
        while (1);
    }

    // XIAO nRF52840 Sense RGB LED (active-low)
    // LEDR / LEDG / LEDB defined by the Seeed board package
    pinMode(LEDR, OUTPUT);
    pinMode(LEDG, OUTPUT);
    pinMode(LEDB, OUTPUT);
    set_rgb(false, false, false);   // all off

    Serial.println("Sleep Stage Classifier Ready.");
}

void loop() {
    // Collect one IMU sample at 1 Hz
    imu_buffer[sample_count][0] = imu.readFloatAccelX();
    imu_buffer[sample_count][1] = imu.readFloatAccelY();
    imu_buffer[sample_count][2] = imu.readFloatAccelZ();
    imu_buffer[sample_count][3] = imu.readFloatGyroX();
    imu_buffer[sample_count][4] = imu.readFloatGyroY();
    imu_buffer[sample_count][5] = imu.readFloatGyroZ();

    sample_count++;

    // After 30 seconds, classify and update LED
    if (sample_count >= 30) {
        float features[NUM_FEATURES];   // NUM_FEATURES = 118
        extract_features(imu_buffer, features);

        int stage = predict_sleep_stage(features);

        print_sleep_stage(stage);
        blink_stage(stage);

        sample_count = 0;   // Reset buffer
    }

    delay(1000);   // 1 Hz IMU sampling
}
\end{ccode}


\section{Arduino Firmware --- Feature Extraction}
\label{app:feature_c}

\begin{ccode}
// Extracts IMU statistics from a 30-sample buffer and writes them
// into the feature vector at the indices defined in features_audio.h.
void extract_imu_features(float buf[30][6], float f[NUM_FEATURES]) {

    // --- Axis-wise statistics (mean, std, max, min) ---
    for (int axis = 0; axis < 6; axis++) {
        f[axis * 4 + 0] = compute_mean(buf, axis);
        f[axis * 4 + 1] = compute_std(buf,  axis);
        f[axis * 4 + 2] = compute_max(buf,  axis);
        f[axis * 4 + 3] = compute_min(buf,  axis);
    }

    // --- Orientation-invariant acceleration magnitude ---
    float acc_mag[30];
    for (int i = 0; i < 30; i++) {
        acc_mag[i] = sqrtf(
            buf[i][0] * buf[i][0] +
            buf[i][1] * buf[i][1] +
            buf[i][2] * buf[i][2]
        );
    }
    f[FEAT_ACC_MAG]     = compute_mean_array(acc_mag, 30);
    f[FEAT_ACC_MAG_STD] = compute_std_array(acc_mag,  30);

    // ... (remaining features follow the same pattern)
}

float compute_mean(float buf[30][6], int axis) {
    float sum = 0.0f;
    for (int i = 0; i < 30; i++) sum += buf[i][axis];
    return sum / 30.0f;
}
\end{ccode}


\section{Arduino Firmware --- LED Feedback}
\label{app:led}

\begin{ccode}
// Helper: sets the three RGB channels simultaneously.
// The XIAO nRF52840 Sense RGB LED is active-low: LOW = on, HIGH = off.
static inline void set_rgb(bool r, bool g, bool b) {
    digitalWrite(LEDR, r ? LOW : HIGH);
    digitalWrite(LEDG, g ? LOW : HIGH);
    digitalWrite(LEDB, b ? LOW : HIGH);
}

// Updates the onboard RGB LED to reflect the classified sleep stage.
// Wake = Red  |  NREM = Green  |  REM = Blue
void blink_stage(int sleep_stage) {
    if      (sleep_stage == 0) set_rgb(true,  false, false);  // Red   -- Wake
    else if (sleep_stage == 1) set_rgb(false, true,  false);  // Green -- NREM
    else                       set_rgb(false, false, true );  // Blue  -- REM

    delay(1500);
    set_rgb(false, false, false);   // off between epochs
}
\end{ccode}


\section{Performance Optimisation}
\label{app:optimisation}

\begin{ccode}
// --- Single-precision float with hardware FPU ---
// The Cortex-M4F executes float32 arithmetic natively in one cycle.
// No fixed-point conversion is required or used.
float acc_mag = sqrtf(ax*ax + ay*ay + az*az);   // FPU sqrtf

// --- Streaming audio accumulator ---
// A single 256-sample PDM window is held in memory at a time.
// Running sums accumulate epoch-level statistics without storing
// the full 30-second waveform (~480 KB at 16 kHz).

typedef struct {
    double energy_sum;     // sum of per-window RMS values
    double energy_sq_sum;  // sum of squared RMS (for std)
    int    window_count;   // number of windows processed
    // ... additional spectral and snoring accumulators
} AudioAccum;

void accum_audio_window(AudioAccum *a, float rms) {
    a->energy_sum    += rms;
    a->energy_sq_sum += rms * rms;
    a->window_count  += 1;
}

// --- Online patient baseline (O(1) per epoch) ---
// Patient-specific normalisation does not require a startup snapshot.
// A running mean is updated after every completed epoch, converging
// to the true individual mean as the session progresses.

pat.acc_sum  += acc_mag;
pat.n_epochs += 1;
double p_acc_mean = pat.acc_sum / pat.n_epochs;   // exact running mean
\end{ccode}


\section{Unit Testing}
\label{app:testing}

\begin{ccode}
// Unit test: verifies feature extraction against a known synthetic input.
// Input: constant acceleration of 1.0g along x-axis, all other axes zero.
// Expected: mean(ax) == 1.0 (FEAT_AX_MEAN, index 0).

void test_feature_extraction() {
    float test_buffer[30][6];

    for (int i = 0; i < 30; i++) {
        test_buffer[i][0] = 1.0f;   // ax = 1g
        test_buffer[i][1] = 0.0f;
        test_buffer[i][2] = 0.0f;
        test_buffer[i][3] = 0.0f;
        test_buffer[i][4] = 0.0f;
        test_buffer[i][5] = 0.0f;
    }

    float features[NUM_FEATURES];
    extract_imu_features(test_buffer, features);

    if (fabsf(features[FEAT_AX_MEAN] - 1.0f) < 0.01f) {
        Serial.println("PASS: mean(ax) == 1.0");
    } else {
        Serial.print("FAIL: mean(ax) = ");
        Serial.println(features[FEAT_AX_MEAN]);
    }
}
\end{ccode}
